\section{Results and Discussion}
\label{sec:results}

Table~\ref{tab:main_results} first gives an overall comparison of the best
results achieved by all methods under the same experimental setting. We
report the best validation accuracy, the epoch at which it is achieved,
and the corresponding test accuracy for each method. The following
subsections then discuss these methods in more detail.

\begin{table*}[t]
\centering
\caption{Overall comparison of the best results for all methods. All methods use the same backbone, dataset, and training budget unless otherwise noted.}
\label{tab:main_results}
\small
\setlength{\tabcolsep}{4pt}   % 默认大约 6pt，可适当缩小
\renewcommand{\arraystretch}{1.05}
\begin{tabular}{p{3.0cm}cccccp{3.6cm}}
\hline
Method & Test Acc & Best Val Acc & Best Epoch & Train Loss @ Best & Val Loss @ Best & Notes \\
\hline
Baseline & 79.65 & 80.42 & 74 & 0.3236 & 0.9317 & Clean reference setting \\
Weight Decay & -- & -- & -- & -- & -- & AdamW with weight decay \\
Dropout & -- & -- & -- & -- & -- & Dropout-enabled fine-tuning \\
Data Augmentation & -- & -- & -- & -- & -- & Stronger augmentation pipeline \\
Label Smoothing + Early Stopping & 83.58 & 83.86 & 92 & 1.0059 & 1.3621 & $\epsilon = 0.1$, patience=25 \\
LLRD + Staged Fine-tuning & 85.33 & 86.06 & 22 & 0.5026 & 0.5633 & Proposed optimization strategy \\
\hline
\end{tabular}
\end{table*}

\subsection{Baseline Results}
\label{sec:baseline}

\paragraph{Reference run.}
We evaluate the baseline method defined in \cref{sec:method:baseline}
under the shared training and evaluation protocol of
\cref{sec:exp:training,sec:exp:eval}. This run serves as the common
reference point for all later comparisons.

\paragraph{Headline numbers.}
The baseline reaches its best validation accuracy of \textbf{80.42\%}
at epoch \textbf{74}, with a validation loss of \textbf{0.9317}. Using
the checkpoint from this epoch, the model achieves a test accuracy of
\textbf{79.65\%}. These results are used as the reference values in the
later method comparisons.

\begin{figure}[!htbp]
  \centering
  % \includegraphics[width=0.85\linewidth]{baseline_curves.png}
  \fbox{\rule{0pt}{4cm}\rule{0.8\linewidth}{0pt}}
  \caption{Baseline ViT-Tiny/16 on CIFAR-100: training and validation
  loss and accuracy across 100 epochs.}
  \label{fig:baseline-curves}
\end{figure}

\paragraph{Discussion.}
After epoch 74, the validation accuracy drops slightly instead of
continuing to improve, which suggests mild overfitting in the later
stage of training. This makes the baseline a reasonable starting point:
it is strong enough to be competitive, but it still leaves room for
later methods to improve generalization.

\subsubsection{Weight Decay Results}
\label{sec:wd}

% Owner: Member 2. Target 1.5--2 pages.

\paragraph{Experiment focus.}
This subsection reports the results of the weight-decay sweep defined in
\cref{sec:method:wd}, evaluated under the shared protocol of
\cref{sec:exp:training,sec:exp:eval}. The only variable that changes is
the AdamW decoupled weight-decay coefficient $\lambda$.

\paragraph{Aggregate numbers.}
\Cref{tab:wd} reports best validation accuracy, final test accuracy,
the train--val gap at the best-val epoch, and wall-clock training
time for each weight-decay setting.

% Weight-decay sweep results. Owner: Member 2.
% Replace -- placeholders with real numbers from the runs.
% Convention: best value per column in \textbf{...}; consistent decimals.

\begin{table}[t]
  \centering
  \small
  \caption{Effect of AdamW weight-decay strength $\lambda$ on
  ViT-Tiny/16 fine-tuned on CIFAR-100. All other hyperparameters fixed
  per \cref{sec:exp:training}. Best per column in bold.}
  \label{tab:wd}
  \begin{tabular}{lcccc}
    \toprule
    Weight decay $\lambda$ & Best val acc (\%) $\uparrow$
                          & Test acc (\%) $\uparrow$
                          & Train--val gap (pp) $\downarrow$
                          & Train time (min) \\
    \midrule
    $0$                   & 84.50 & \textbf{84.75} & 5.38 & 31.5 \\
    $10^{-4}$             & 84.38 & 84.31 & 5.28 & 31.5 \\
    $10^{-2}$             & 84.28 & 84.74 & 5.54 & \textbf{31.3} \\
    $5 \times 10^{-2}$    & 83.88 & 84.05 & 5.54 & 31.5 \\
    $10^{-1}$             & \textbf{84.56} & 84.33 & \textbf{4.97} & 31.7 \\
    \bottomrule
  \end{tabular}
\end{table}


\paragraph{Learning curves.}
\Cref{fig:wd-curves} overlays training and validation loss and
accuracy across all five settings, sharing axes with
\cref{fig:baseline-curves} for direct comparison.

\begin{figure}[t]
  \centering
  % \includegraphics[width=0.9\linewidth]{wd_curves.png}  % uncomment when figure is ready
  \fbox{\rule{0pt}{4cm}\rule{0.85\linewidth}{0pt}}
  \caption{Training and validation curves across weight-decay
  strengths $\lambda \in \{0, 10^{-4}, 10^{-2}, 5\times 10^{-2},
  10^{-1}\}$. \draftnote{Drop \texttt{figures/wd\_curves.png} and
  replace the \texttt{\textbackslash fbox} placeholder with the
  commented \texttt{\textbackslash includegraphics} above. One
  figure, four subplots (train loss, val loss, train acc, val acc),
  one line per $\lambda$, color-coded.}}
  \label{fig:wd-curves}
\end{figure}

\paragraph{Train--val gap.}
\Cref{fig:wd-gap} plots the train--val accuracy gap as a function of
training epoch for each $\lambda$. A flat or shrinking gap at large
$\lambda$ is direct evidence that weight decay is constraining
memorization rather than just trading capacity for regularization.

\begin{figure}[t]
  \centering
  % \includegraphics[width=0.7\linewidth]{wd_gap.png}  % uncomment when figure is ready
  \fbox{\rule{0pt}{4cm}\rule{0.65\linewidth}{0pt}}
  \caption{Train--val accuracy gap over training for each weight-decay
  strength. \draftnote{Drop \texttt{figures/wd\_gap.png} and swap
  in the commented \texttt{\textbackslash includegraphics} above. The
  caption should state the punch line: e.g. ``$\lambda=10^{-2}$
  reduces the final-epoch gap by X percentage points relative to
  $\lambda=0$.''}}
  \label{fig:wd-gap}
\end{figure}

\paragraph{Best validation accuracy as a function of $\lambda$.}
\Cref{fig:wd-best-val} summarizes the sweep by plotting best
validation accuracy against $\lambda$ on a log-x axis. This is the
single plot that answers ``what coefficient should I pick?''

\begin{figure}[t]
  \centering
  % \includegraphics[width=0.6\linewidth]{wd_best_val_vs_wd.png}  % uncomment when ready
  \fbox{\rule{0pt}{4cm}\rule{0.55\linewidth}{0pt}}
  \caption{Best validation accuracy versus weight-decay strength
  $\lambda$ (log scale on the x-axis; $\lambda=0$ is shown as a
  separate marker on the left). \draftnote{Drop
  \texttt{figures/wd\_best\_val\_vs\_wd.png} and swap in the commented
  \texttt{\textbackslash includegraphics} above. Mark the optimum.}}
  \label{fig:wd-best-val}
\end{figure}

\paragraph{Is weight decay effective for fine-tuning ViT-Tiny on
CIFAR-100?}
\label{sec:wd:discussion}
\draftnote{Topic sentence: state whether the best $\lambda > 0$
beats $\lambda = 0$ on test accuracy, and by how many points.
Then 2--3 sentences of evidence: cite the relevant rows of
\cref{tab:wd} and the corresponding curve in \cref{fig:wd-curves}.}

\paragraph{What is roughly the optimal coefficient?}
\draftnote{Topic sentence: name the optimum (likely $10^{-2}$ or
$5\times 10^{-2}$ for AdamW on a small ViT). Then describe the shape
of the curve in \cref{fig:wd-best-val}: is it sharply peaked or a
broad plateau? A broad plateau is the more useful finding because it
implies the practitioner does not need fine tuning.}

\paragraph{Does it act as overfitting prevention or capacity
suppression?}
\draftnote{Use the train--val gap evidence in \cref{fig:wd-gap}.
The decision rule: if increasing $\lambda$ \emph{shrinks} the gap
\emph{while} val accuracy improves or holds, weight decay is
preventing overfitting. If the gap shrinks because train accuracy
collapses faster than val, then weight decay is suppressing capacity
(underfitting). For our largest $\lambda = 10^{-1}$, characterize
which of these two regimes we are in; this is the key qualitative
finding of the section.}

\paragraph{Limitations.}
\label{sec:wd:limitations}
\draftnote{1 short paragraph. Cover: (i) single seed per
configuration, so reported differences may sit within run-to-run
noise---report the magnitude if multi-seed runs are available;
(ii) we study fine-tuning from an ImageNet-pretrained checkpoint,
not training from scratch, so conclusions may not transfer to the
from-scratch regime where weight decay matters more
\citep{touvron2021deit}; (iii) absolute differences between
neighboring $\lambda$ values may be small (sub-1\%), so we are
careful not to over-interpret rank orderings.}

\subsection{Dropout Results}
\label{sec:dropout}

% Owner: Member 3. Bullet-level skeleton; expand to ~1 page.

\paragraph{Experiment focus.}
We evaluate the dropout variants defined in
\cref{sec:method:dropout} under the shared protocol of
\cref{sec:exp:training,sec:exp:eval}. Only the dropout placement and
rate vary across runs. Overall, the results from the 30-epoch training run generally outperformed those from the 100-epoch run, indicating that longer training leads to overfitting; therefore, we primarily analyze the generalization performance under the 100-epoch setting.

\paragraph{Results summary.} For standard Dropout, as the drop rate increases, the train-val gap decreases significantly (10.67 → 8.98 pp), but test accuracy peaks at a drop rate of 0.3 (82.90\%), indicating that overly strong regularization does not necessarily lead to better test performance. In contrast, DropPath demonstrates the most stable performance, achieving both the highest test accuracy (83.56\%) and the lowest gap (4.64 pp) at 100 epochs. R-Drop yields a modest improvement (82.94\%) through consistency constraints. Attention Dropout and Patch Dropout show limited gains, with best test accuracies of 82.34\% and 82.29\%, respectively.


\begin{table}[t]
\centering
\small
\caption{Effect of dropout rate. The 100-epoch results are the main setting, while 30-epoch results serve as a short-horizon reference.}
\label{tab:dropout}
\resizebox{\columnwidth}{!}{
\begin{tabular}{lccccc}
\toprule
Drop rate
& \multicolumn{3}{c}{100 epochs}
& \multicolumn{2}{c}{30 epochs} \\
\cmidrule(lr){2-4}\cmidrule(lr){5-6}
      & Best val (\%) $\uparrow$
      & Test (\%) $\uparrow$
      & Gap (pp) $\downarrow$
      & Best val (\%) $\uparrow$
      & Test (\%) $\uparrow$ \\
\midrule
0.1 & 82.08 & 82.39 & 10.67 & 84.44 & 84.26 \\
0.2 & 82.06 & 82.49 & 10.70 & 84.50 & \textbf{84.74} \\
0.3 & 82.10 & \textbf{82.90} & 10.29 & 84.36 & 84.67 \\
0.5 & \textbf{82.88} & 82.87 & \textbf{8.98} & \textbf{84.68} & 84.50 \\
\bottomrule
\end{tabular}
}
\end{table}

\begin{table}[t]
\centering
\small
\caption{Effect of DropPath rate.}
\label{tab:droppath}
\resizebox{\columnwidth}{!}{
\begin{tabular}{lccccc}
\toprule
DropPath
& \multicolumn{3}{c}{100 epochs}
& \multicolumn{2}{c}{30 epochs} \\
\cmidrule(lr){2-4}\cmidrule(lr){5-6}
      & Best val (\%) $\uparrow$
      & Test (\%) $\uparrow$
      & Gap (pp) $\downarrow$
      & Best val (\%) $\uparrow$
      & Test (\%) $\uparrow$ \\
\midrule
0.05 & 83.30 & 82.81 & 9.13  & 84.86 & \textbf{85.05} \\
0.1  & 83.56 & 83.51 & 7.86  & \textbf{85.42} & 84.94 \\
0.2  & 83.62 & 83.17 & 6.55  & 85.06 & 84.61 \\
0.3  & \textbf{83.68} & \textbf{83.56} & \textbf{4.64} & 83.80 & 83.51 \\
\bottomrule
\end{tabular}
}
\end{table}

\begin{table}[t]
\centering
\small
\caption{Effect of R-Drop.}
\label{tab:rdrop}
\resizebox{\columnwidth}{!}{
\begin{tabular}{lccccc}
\toprule
Drop rate
& \multicolumn{3}{c}{100 epochs}
& \multicolumn{2}{c}{30 epochs} \\
\cmidrule(lr){2-4}\cmidrule(lr){5-6}
      & Best val (\%) $\uparrow$
      & Test (\%) $\uparrow$
      & Gap (pp) $\downarrow$
      & Best val (\%) $\uparrow$
      & Test (\%) $\uparrow$ \\
\midrule
0.1 & 82.36 & 82.68 & 10.88 & 85.02 & 84.83 \\
0.2 & \textbf{83.12} & 82.55 & \textbf{9.68} & 84.76 & \textbf{84.84} \\
0.3 & 82.56 & \textbf{82.94} & 10.21 & \textbf{85.28} & 84.62 \\
\bottomrule
\end{tabular}
}
\end{table}

\begin{table}[t]
\centering
\small
\caption{Effect of attention dropout rate.}
\label{tab:dropattn}
\resizebox{\columnwidth}{!}{
\begin{tabular}{lccccc}
\toprule
Drop rate
& \multicolumn{3}{c}{100 epochs}
& \multicolumn{2}{c}{30 epochs} \\
\cmidrule(lr){2-4}\cmidrule(lr){5-6}
      & Best val (\%) $\uparrow$
      & Test (\%) $\uparrow$
      & Gap (pp) $\downarrow$
      & Best val (\%) $\uparrow$
      & Test (\%) $\uparrow$ \\
\midrule
0.05 & 81.90 & 81.86 & 11.67 & 84.12 & 84.45 \\
0.1  & \textbf{82.50} & \textbf{82.34} & \textbf{11.02} & 84.16 & 84.54 \\
0.2  & 81.46 & 82.03 & 11.85 & 84.16 & \textbf{84.55} \\
0.3  & 81.82 & 82.03 & 11.20 & \textbf{84.70} & 84.06 \\
\bottomrule
\end{tabular}
}
\end{table}

\begin{table}[t]
\centering
\small
\caption{Effect of patch dropout rate.}
\label{tab:droppatch}
\resizebox{\columnwidth}{!}{
\begin{tabular}{lccccc}
\toprule
Drop rate
& \multicolumn{3}{c}{100 epochs}
& \multicolumn{2}{c}{30 epochs} \\
\cmidrule(lr){2-4}\cmidrule(lr){5-6}
  & Best val (\%) $\uparrow$
      & Test (\%) $\uparrow$
      & Gap (pp) $\downarrow$
      & Best val (\%) $\uparrow$
      & Test (\%) $\uparrow$ \\
\midrule
0.1 & 81.98 & 81.74 & 11.48 & \textbf{84.44} & \textbf{84.70} \\
0.2 & \textbf{82.24} & \textbf{82.29} & \textbf{10.73} & 84.16 & 84.24 \\
0.3 & 82.10 & 81.61 & 11.00 & 84.10 & 83.91 \\
\bottomrule
\end{tabular}
}
\end{table}

\subsubsection{Data Augmentation Results}
\label{sec:augmentation}

% Owner: Member 4. Bullet-level skeleton; expand to ~1--1.5 pages
% (this section can be the most visual one).

\paragraph{Experiment focus.}
We evaluate the augmentation families defined in \cref{sec:method:aug}
under the shared protocol of \cref{sec:exp:training,sec:exp:eval}.
Only the data pipeline changes; the model, optimizer, schedule, and
checkpoint-selection rule remain fixed.

\paragraph{Results.}
\draftnote{Expand this subsection into prose using the checklist below.}
\begin{itemize}
  \item Table: per augmentation (or per combination), best val acc /
        test acc / train--val gap / train time.
  \item Curves figure overlaying val accuracy per augmentation.
  \item Optional: a robustness probe (e.g. corruption noise) if time
        permits; otherwise skip.
\end{itemize}

\paragraph{Discussion.}
\label{sec:aug:discussion}
\draftnote{Answer the three Member 4 questions in short paragraphs.}
\begin{enumerate}
  \item Which augmentation is most effective?
  \item How big is the gap between simple and strong augmentation?
  \item Is data-level regularization more effective than
        parameter-level (cross-reference \cref{sec:wd,sec:dropout})?
\end{enumerate}

\subsubsection{Label Smoothing and Early Stopping Results}
\label{sec:lsearly}

% Owner: Member 5. Bullet-level skeleton; expand to ~1 page.

\paragraph{Label smoothing results.}
\draftnote{Expand this subsection into prose using the checklist below.}
\begin{itemize}
  \item Reference the method definition in \cref{sec:method:lsearly}
        and note that all shared training details follow
        \cref{sec:exp:training,sec:exp:eval}.
  \item Report best val acc / test acc / train--val gap / train time
        in a table.
  \item Briefly discuss confidence calibration (e.g. softmax max
        probability or ECE), since this is where label smoothing
        \citep{szegedy2016rethinking,muller2019labelsmoothing}
        differs most from other regularizers.
\end{itemize}

\paragraph{Early stopping results.}
\draftnote{Expand this subsection into prose using the checklist below.}
\begin{itemize}
  \item Reference the method definition in \cref{sec:method:lsearly}
        and note the shared protocol in
        \cref{sec:exp:training,sec:exp:eval}.
  \item Compare against fixed 30-epoch training: at which epoch does
        early stopping trigger? How much compute saved?
  \item Test accuracy at the early-stopped checkpoint vs final-epoch
        checkpoint.
\end{itemize}

\paragraph{Discussion.}
\label{sec:lsearly:discussion}
\draftnote{Answer the three Member 5 questions in short paragraphs.}
\begin{enumerate}
  \item Does label smoothing improve generalization here?
  \item Does early stopping actually help, or is it just ``stop
        early''?
  \item How do these output- and process-level regularizers relate to
        the parameter-level (\cref{sec:wd,sec:dropout}) and
        data-level (\cref{sec:augmentation}) ones?
\end{enumerate}





\subsection{LLRD + Staged Fine-tuning Results}
\label{sec:llrd}

We next evaluate the proposed optimization strategy based on layer-wise
learning rate decay (LLRD) and staged fine-tuning. Under the same
backbone, dataset, and training budget as the baseline, this method
achieves a best validation accuracy of \textbf{86.06\%} at epoch
\textbf{22}, with a corresponding test accuracy of \textbf{85.33\%}.
Compared with the baseline result of 81.86\% validation accuracy, this
gives an improvement of \textbf{4.20} percentage points.

The training curves show that this method improves much faster after the
model is switched from head-only training to full-model fine-tuning.
Validation accuracy rises quickly in the early part of Stage 2, then
reaches its best result around epoch 22. After that point, training
accuracy continues to increase while validation accuracy stops
improving, which suggests the start of overfitting. This result shows
that LLRD with staged fine-tuning is an effective and stable way to
adapt a pretrained ViT to CIFAR-100, especially under a limited
training budget.